\section*{Hoofdstuk \ref{ch:g7:atoms-and-molecules} --- Atomen en moleculen}

\begin{solution}{exo:g7:atoms-and-molecules:1}
Een \omterm{def:g7:atoms-and-molecules:atom}{atoom} is een van de uiterst kleine deeltjes waaruit alle materie is
opgebouwd (er bestaan ongeveer honderd soorten). Een \omterm{def:g7:atoms-and-molecules:molecule}{molecuul} is een
groep \omterm{def:g7:atoms-and-molecules:atom}{atomen} die stevig aan elkaar vastzitten, het kleinste stukje van
een stof dat nog die stof is.
\end{solution}

\begin{solution}{exo:g7:atoms-and-molecules:2}
\ce{H}, \ce{C}, \ce{N}, \ce{O}, \ce{Na}, \ce{Fe}.
\end{solution}

\begin{solution}{exo:g7:atoms-and-molecules:3}
1 koolstofatoom en 2 zuurstofatomen: in totaal 3 \omterm{def:g7:atoms-and-molecules:atom}{atomen}.
\end{solution}

\begin{solution}{exo:g7:atoms-and-molecules:4}
Water, zuurstof, methaan, stikstof.
\end{solution}

\begin{solution}{exo:g7:atoms-and-molecules:5}
\ce{NO2}.
\end{solution}

\begin{solution}{exo:g7:atoms-and-molecules:6}
\ce{O2} is één \omterm{def:g7:atoms-and-molecules:molecule}{molecuul} uit twee aan elkaar vastzittende zuurstofatomen;
\ce{2O} zijn twee losse zuurstofatomen. \ce{CO} is een \omterm{def:g7:atoms-and-molecules:molecule}{molecuul} uit één
koolstofatoom en één zuurstofatoom (twee hoofdletters, twee symbolen);
\ce{Co} is één \omterm{def:g7:atoms-and-molecules:atom}{atoom} kobalt (één \omterm{def:g7:atoms-and-molecules:symbol}{symbool}, een hoofdletter en dan een
kleine letter).
\end{solution}

\begin{solution}{exo:g7:atoms-and-molecules:7}
\ce{3H2O}: $3\times 2 = 6$ waterstofatomen en $3\times 1 = 3$
zuurstofatomen. \ce{5O2}: geen waterstof, $5\times 2 = 10$
zuurstofatomen.
\end{solution}

\begin{solution}{exo:g7:atoms-and-molecules:8}
Zwart is koolstof, rood is zuurstof: één koolstof, twee zuurstof,
\ce{CO2}, koolstofdioxide.
\end{solution}

\begin{solution}{exo:g7:atoms-and-molecules:9}
Het water van Dalton was \ce{HO}. Het echte \omterm{def:g7:atoms-and-molecules:molecule}{molecuul}, \ce{H2O}, heeft
één waterstofatoom meer: het \omterm{def:g7:atoms-and-molecules:molecule}{molecuul} van Dalton kwam één \omterm{def:g7:atoms-and-molecules:atom}{atoom} tekort.
\end{solution}

\begin{solution}{exo:g7:atoms-and-molecules:10}
Zuiver water is geen \omterm{def:g6:pure-substances-and-mixtures:mixture}{mengsel}: al zijn \omterm{def:g7:atoms-and-molecules:molecule}{moleculen} zijn gelijke
\ce{H2O}-moleculen. \omterm{def:g4:air-a-mixture-of-gases:air}{Lucht} is een \omterm{def:g6:pure-substances-and-mixtures:mixture}{mengsel}: het bevat stikstofmoleculen,
zuurstofmoleculen, argonatomen en een paar koolstofdioxidemoleculen.
\end{solution}

\begin{solution}{exo:g7:atoms-and-molecules:11}
Tien miljoen \omterm{def:g7:atoms-and-molecules:atom}{atomen} maken \qty{1}{mm}, dus voor \qty{1}{cm} $= \qty{10}{mm}$
zijn $10 \times$ tien miljoen $=$ honderd miljoen \omterm{def:g7:atoms-and-molecules:atom}{atomen} nodig.
$\qty{100}{m} = \num{100000}$~mm vraagt $\num{100000}\times$ tien miljoen
$=$ duizend miljard \omterm{def:g7:atoms-and-molecules:atom}{atomen}.
\end{solution}

\begin{solution}{exo:g7:atoms-and-molecules:12}
$6 + 12 + 6 = 24$ \omterm{def:g7:atoms-and-molecules:atom}{atomen}. Waterstof: $\frac{12}{24} = \qty{50}{\percent}$.
Koolstof: $\frac{6}{24} = \qty{25}{\percent}$ (en zuurstof de overige
\qty{25}{\percent}).
\end{solution}

\begin{solution}{pb:g7:atoms-and-molecules:1}
\textbf{1.} Stikstof \ce{N2}, 2 \omterm{def:g7:atoms-and-molecules:atom}{atomen}; zuurstof \ce{O2}, 2 \omterm{def:g7:atoms-and-molecules:atom}{atomen};
koolstofdioxide \ce{CO2}, 3 \omterm{def:g7:atoms-and-molecules:atom}{atomen}.

\textbf{2.} Argonatomen vormen geen \omterm{def:g7:atoms-and-molecules:molecule}{moleculen}: argon bestaat uit losse
\omterm{def:g7:atoms-and-molecules:atom}{atomen}.

\textbf{3.} Stikstof: $78 \times 2 = 156$ \omterm{def:g7:atoms-and-molecules:atom}{atomen}. Zuurstof: $21\times 2 =
42$ \omterm{def:g7:atoms-and-molecules:atom}{atomen}.

\textbf{4.} $156 + 42 + 1 = 199$ \omterm{def:g7:atoms-and-molecules:atom}{atomen}.

\textbf{5.} $\frac{156}{199} \approx 0.78$, ongeveer \qty{78}{\percent}.

\textbf{6.} In zuurstofmoleculen: $16\times 2 = 32$; in
koolstofdioxidemoleculen: $4 \times 2 = 8$; in totaal $32 + 8 = 40$.

\textbf{7.} Er ontbreken $42 - 40 = 2$ zuurstofatomen: één
zuurstofmolecuul. Ze zijn terechtgekomen in watermoleculen, \ce{H2O}, die
het lichaam maakt met de waterstof uit zijn voedsel; dat water verlaat
het lichaam deels als waterdamp in de adem (die hier niet meetelt).

\textbf{8.} Ingeademd: 0. Uitgeademd: 4 (één in elk \ce{CO2}).

\textbf{9.} Uit het voedsel dat het lichaam gebruikt: suikers en vetten
bevatten koolstofatomen, en het lichaam geeft die af in het
koolstofdioxide dat het uitademt.

\textbf{10.} Er zijn $21 - 16 = 5$ zuurstofmoleculen verdwenen en 4
koolstofdioxidemoleculen bijgekomen.

\textbf{11.} Stikstof: twee blauwe balletjes, verbonden door drie
staafjes. Zuurstof: twee rode balletjes, verbonden door twee staafjes.
Koolstofdioxide: een zwart balletje tussen twee rode balletjes, elk met
twee staafjes eraan vast, op een rechte lijn. Argon: één los lichtblauw
balletje.

\textbf{12.} $156 + 32 + 12 + 1 = 201$ \omterm{def:g7:atoms-and-molecules:atom}{atomen} (de 4
koolstofdioxidemoleculen bevatten 12 \omterm{def:g7:atoms-and-molecules:atom}{atomen}). Dat is $201 - 199 = 2$
\omterm{def:g7:atoms-and-molecules:atom}{atomen} meer dan in de ingeademde \omterm{def:g4:air-a-mixture-of-gases:air}{lucht}: 4 koolstofatomen gewonnen uit
voedsel, min de 2 zuurstofatomen die als water zijn vertrokken.
\end{solution}
